\section{弹性体、物质坐标与小形变}

弹性（elasticity）指撤去使物体形变的作用后，物体能够恢复原有构形的性质；
“完全弹性”不等于“任意大形变都服从线性胡克定律”。固体可以储存剪切弹性能，
流体则主要以黏性应力抵抗形变速率。黏弹性材料的应力还与加载历史有关；
塑性形变可以留下不可恢复的部分。非牛顿流体、黏弹性与塑性属于流变学
（rheology）研究的范围，不能把三者简单等同。本章先研究小形变、可逆、线性弹性体。

用参考构形中某材料点的位置 $\vect R=(X,Y,Z)$ 给它作永久标记。
该点在时刻 $t$ 的实际位置为 $\vect r(\vect R,t)$，位移为
\begin{equation}
 \vect u(\vect R,t)=\vect r(\vect R,t)-\vect R.
 \label{ch12:eq:displacement}
\end{equation}
$\vect R$ 是“哪一个材料点”，$\vect r$ 是“它现在在哪里”；不能把位移场误认为速度场。
若相邻点的原间隔为 $\dd\vect R$，泰勒展开到一阶给出
\begin{equation}
 \dd u_i=\sum_{j=1}^{3}\frac{\partial u_i}{\partial X_j}\dd X_j,\qquad
 \dd\vect r=(\mathsf I+\mathsf H)\dd\vect R,\quad H_{ij}=\frac{\partial u_i}{\partial X_j}.
 \label{ch12:eq:gradient}
\end{equation}
这里 $X_1=X,X_2=Y,X_3=Z$，$\mathsf H$ 是九个分量组成的线性映射；
本章采用“行是位移分量、列是求导坐标”的矩阵约定。整体平移改变 $\vect u$，
却不改变相邻点的间隔，因而不能算作形变。

任何矩阵都有唯一分解
\[
 \mathsf H=\frac{\mathsf H+\mathsf H^{\mathsf T}}2+
            \frac{\mathsf H-\mathsf H^{\mathsf T}}2
          \equiv\boldsymbol\varepsilon+\mathsf W.
\]
上标 $\mathsf T$ 表示转置。$\boldsymbol\varepsilon$ 对称，$\mathsf W$ 反对称。
例如刚体发生小转角 $\delta\vect\phi$ 时，
$\dd\vect u=\delta\vect\phi\times\dd\vect R$，相应矩阵为
\[
 \mathsf W=
 \begin{pmatrix}
 0&-\delta\phi_z&\delta\phi_y\\
 \delta\phi_z&0&-\delta\phi_x\\
 -\delta\phi_y&\delta\phi_x&0
 \end{pmatrix}.
\]
它只转动微元，不改变微元各边长度和夹角。反过来，由
\[
 |\dd\vect r|^2-|\dd\vect R|^2
 =2\dd\vect R\cdot\boldsymbol\varepsilon\,\dd\vect R+O(H^2)
\]
可见，一阶长度变化只取决于对称部分。因此定义小应变张量（strain tensor）
\begin{equation}
 \varepsilon_{ij}=\frac12\left(\frac{\partial u_i}{\partial X_j}
                           +\frac{\partial u_j}{\partial X_i}\right),\qquad
 \boldsymbol\varepsilon=
 \begin{pmatrix}
 \varepsilon_{xx}&\gamma_{xy}/2&\gamma_{xz}/2\\
 \gamma_{xy}/2&\varepsilon_{yy}&\gamma_{yz}/2\\
 \gamma_{xz}/2&\gamma_{yz}/2&\varepsilon_{zz}
 \end{pmatrix}.
 \label{ch12:eq:strain}
\end{equation}
这一分解适用于小位移梯度；大角度刚体转动不能直接用截断的一阶应变判断。

\subsection{正应变与剪应变的几何意义}

沿 $X$ 方向原长为 $L$ 的小线段，伸长为 $\Delta L$，正应变为
$\varepsilon_{xx}=\Delta L/L=\partial u_x/\partial X$。
正号表示伸长，负号表示缩短。两个原来正交的单位线段，变形后方向分别近似为
$\uvec x+\partial_X\vect u$ 和 $\uvec y+\partial_Y\vect u$。
它们点积到一阶为 $\partial_Yu_x+\partial_Xu_y$。
若夹角从 $\pi/2$ 变成 $\pi/2-\alpha$，因
$\cos(\pi/2-\alpha)=\sin\alpha\simeq\alpha$，故
\begin{equation}
 \gamma_{xy}=2\varepsilon_{xy}
 =\frac{\partial u_x}{\partial Y}+\frac{\partial u_y}{\partial X}
 \simeq\alpha.
 \label{ch12:eq:shear-angle}
\end{equation}
$\gamma_{xy}$ 是工程剪应变，不能把它与张量分量 $\varepsilon_{xy}$ 混用。
例如 $u_x=\gamma Y,u_y=0$ 时，上边相对下边横移 $\gamma H$，
$\tan\alpha=\gamma\simeq\alpha$。正应变、剪应变均无量纲。
图\ref{ch12:fig:strain}还说明：同一受拉材料内，微元取向不同，看到的正应变、剪应变分量也不同。

\begin{center}
\begin{tikzpicture}[>=Latex,scale=.87]
 \draw[guide] (0,0) rectangle(4,1.25);
 \draw[thick] (-.35,.12) rectangle(4.35,1.13);
 \draw[vector] (-.35,.63)--(-1.3,.63);
 \draw[vector] (4.35,.63)--(5.3,.63);
 \draw[thick,DeepBlue] (.5,.4) rectangle(1.3,.85);
 \draw[thick,DeepBlue] (2.5,.25)--(3.15,.63)--(2.5,1)--(1.85,.63)--cycle;
 \node at (2,-.55) {同一受拉杆内的不同取向微元};
 \begin{scope}[shift={(7.1,0)}]
  \draw[guide] (0,0) rectangle(1.3,1.25);
  \draw[thick,DeepBlue] (0,0) rectangle(1.8,1.05);
  \draw[<->] (0,-.4)--(1.3,-.4) node[midway,below] {$L$};
  \node at (.9,1.75) {正应变};
 \end{scope}
 \begin{scope}[shift={(10.5,0)}]
  \draw[guide] (0,0) rectangle(1.3,1.25);
  \draw[thick,DeepBlue] (0,0)--(1.3,0)--(1.8,1.25)--(.5,1.25)--cycle;
  \draw[<->] (0,1.65)--(.5,1.65);
  \node[above] at (.25,1.65) {$\gamma H$};
  \draw (0,.65) arc(90:68.20:.65);
  \node[left] at (-.12,.55) {$\alpha$};
  \node at (.9,-.8) {剪应变};
 \end{scope}
\end{tikzpicture}
\figcap{受拉材料中的正应变与剪应变。虚线为参考构形；右侧将微元平移、转动对齐后比较长度和直角的变化。}
\label{ch12:fig:strain}
\end{center}

\section{有向截面、应力与弹性定律}

\subsection{为什么内部受力需要应力张量}

只指定物体内部一点，尚不能指定“谁对谁施力”。需在该点切出法向为 $\vect n$ 的小面，
取 $\vect n$ 指向所研究部分的外侧，以外部对内部的面力除以面积，定义面力密度
$\vect t(\vect n)=\dd\vect F/\dd S$，单位为 $\mathrm{N/m^2}$，即帕斯卡。
同一点不同方向的截面通常有不同的面力；所以内部应力（stress）不是一个矢量。

面力与有向面积是线性关系，可以用小四面体证明：斜面的面积为 $\dd S$，
其三个坐标面投影面积分别为 $n_j\dd S$。体积力及惯性力是微元边长的三次量，
面力是二次量；微元趋于零时，主阶平衡给出
$\vect t(\vect n)=\sum_j n_j\vect t(\uvec e_j)$。
负的 $n_j$ 由反向截面的作用反作用关系 $\vect t(-\vect n)=-\vect t(\vect n)$ 处理。
因此定义
\begin{equation}
 t_i(\vect n)=\sum_j\sigma_{ij}n_j,\qquad
 \dd\vect F=\boldsymbol\sigma\,\vect n\,\dd S.
 \label{ch12:eq:traction}
\end{equation}
$\sigma_{ij}$ 表示法向沿第 $j$ 轴的面上，第 $i$ 方向的面力分量。
$\sigma_{nn}=\vect n\cdot\boldsymbol\sigma\vect n$ 为正应力，余下的面内分量为剪应力，
如图\ref{ch12:fig:stress}所示。拉伸正应力取正；流体压强 $p>0$ 对应 $\boldsymbol\sigma=-p\mathsf I$。

\begin{center}
\begin{tikzpicture}[>=Latex,scale=.93]
 \begin{scope}
  \draw[thick] (0,0)--(2.5,0)--(3.25,.8)--(.75,.8)--cycle;
  \coordinate(P) at(1.6,.4);
  \draw[vector] (P)--++(0,1.5) node[above] {$\vect t=\sigma_{nn}\vect n$};
  \draw[->,TextGray] (.55,.32)--(.55,1.15) node[left] {$\vect n$};
  \node at (1.6,-.6) {正应力：垂直截面};
 \end{scope}
 \begin{scope}[shift={(6,0)}]
  \draw[thick] (0,0)--(2.5,0)--(3.25,.8)--(.75,.8)--cycle;
  \coordinate(Q) at(1.35,.35);
  \draw[vector] (Q)--++(1.9,0) node[right] {$\vect t_\parallel$};
  \draw[->,TextGray] (Q)--++(0,1.2) node[above] {$\vect n$};
  \node at (1.6,-.6) {剪应力：平行截面};
 \end{scope}
\end{tikzpicture}
\figcap{正应力与剪应力是同一有向截面面力的垂直、平行分量。箭头起于截面内部，面力单位为力除以面积。}
\label{ch12:fig:stress}
\end{center}

在没有微观偶力矩的普通连续介质中，应力还满足 $\sigma_{ij}=\sigma_{ji}$。
例如考察尺寸 $\dd x,\dd y,\dd z$ 的小长方体绕 $z$ 轴的力矩：
$x$ 法向两面上的 $y$ 向剪力合力矩为 $\sigma_{yx}\dd x\dd y\dd z$，
$y$ 法向两面上的 $x$ 向剪力合力矩为 $-\sigma_{xy}\dd x\dd y\dd z$。
正应力不贡献这一主阶力矩。微元转动惯量比体积多两个长度幂次，
故角动量方程的最低阶要求 $\sigma_{yx}-\sigma_{xy}=0$，其余两对同理。

\subsection{胡克定律、各向同性与独立弹性常数}

在指定坐标系中，应力映射也可显式写成
\[
 \begin{pmatrix}\dd F_x\\\dd F_y\\\dd F_z\end{pmatrix}
 =
 \begin{pmatrix}
 \sigma_{xx}&\sigma_{xy}&\sigma_{xz}\\
 \sigma_{yx}&\sigma_{yy}&\sigma_{yz}\\
 \sigma_{zx}&\sigma_{zy}&\sigma_{zz}
 \end{pmatrix}
 \begin{pmatrix}n_x\\n_y\\n_z\end{pmatrix}\dd S.
\]
对法向沿 $+y$ 的截面，法向分量为 $\sigma_{yy}\dd S$，
切向分量为 $\sigma_{xy}\dd S$、$\sigma_{zy}\dd S$，不能漏掉三维中的第二个切向方向。
小形变范围内的一般线性本构关系写成
$\sigma_{ij}=\sum_{k,l}C_{ijkl}\varepsilon_{kl}$。
四个指标各有三个取值，形式上有 $3^4=81$ 项。
应力和应变分别对称，使独立分量由 $9\times9$ 降为 $6\times6=36$。
若形变完全可逆、存在弹性能密度 $w(\boldsymbol\varepsilon)$，则
$\sigma_{ij}=\partial w/\partial\varepsilon_{ij}$（非对角独立分量要一致计数），
混合二阶导数相等给出 $C_{ijkl}=C_{klij}$，因此最多有 $6\times7/2=21$ 个独立常数。
这是弹性对称性和能量条件共同给出的结论，并非所有材料一开始只有两个常数。

各向同性意味着旋转坐标后本构关系不变。各向同性四阶张量只能由克罗内克符号
$\delta_{ij}$ 的三种两两配对组成；再用前述对称性，便只剩两个常数：
\begin{equation}
 \sigma_{ij}=\lambda_e\,\varepsilon_{kk}\delta_{ij}+2G\varepsilon_{ij},
 \qquad \varepsilon_{kk}=\varepsilon_{xx}+\varepsilon_{yy}+\varepsilon_{zz}.
 \label{ch12:eq:lame}
\end{equation}
这里重复指标 $k$ 表示对三个方向求和；$\delta_{ij}$ 在 $i=j$ 时为1，其余为0。
$\lambda_e,G$ 称为拉梅常数（Lamé constants）；用 $\lambda_e$ 的下标区别后文波长 $\lambda$。
在线性范围，$\varepsilon_{kk}=\Delta V/V$，因为
$(1+\varepsilon_{xx})(1+\varepsilon_{yy})(1+\varepsilon_{zz})-1$
的一阶项正是三者之和。

图\ref{ch12:fig:moduli}给出三种实验操作。纯剪切中
$\sigma_{xy}=2G\varepsilon_{xy}=G\gamma_{xy}$，故 $G$ 就是剪切模量。
静水压缩时 $\sigma_{xx}=\sigma_{yy}=\sigma_{zz}=-p$，取式\eqref{ch12:eq:lame}的迹得到
\begin{equation}
 -3p=(3\lambda_e+2G)\frac{\Delta V}{V},\qquad
 K=-V\frac{\dd p}{\dd V}=\lambda_e+\frac{2G}{3}.
 \label{ch12:eq:bulk}
\end{equation}
压强增加时体积减小，负号使稳定材料的体积模量 $K$ 为正。

\begin{center}
\begin{tikzpicture}[>=Latex,scale=.87]
 \draw[guide] (0,0) rectangle(2,1.5);
 \draw[thick,DeepBlue] (0,0)--(2,0)--(2.5,1.5)--(.5,1.5)--cycle;
 \draw[vector] (.7,1.9)--(2,1.9) node[right] {$\tau$};
 \draw[vector] (1.4,-.4)--(.1,-.4);
 \node at (1.2,-1) {剪切：$\tau=G\gamma$};
 \begin{scope}[shift={(5,0)}]
  \draw[guide] (0,0) rectangle(2,1.5);
  \draw[thick,DeepBlue] (.2,.15) rectangle(1.8,1.35);
  \draw[vector] (1,2.1)--(1,1.4);
  \draw[vector] (1,-.6)--(1,.1);
  \draw[vector] (-.7,.75)--(.15,.75);
  \draw[vector] (2.7,.75)--(1.85,.75);
  \node at (1,-1) {体压缩：$K$};
 \end{scope}
 \begin{scope}[shift={(10,0)}]
  \draw[guide] (0,0) rectangle(2,1.5);
  \draw[thick,DeepBlue] (-.3,.2) rectangle(2.3,1.3);
  \draw[vector] (-.3,.75)--(-1,.75);
  \draw[vector] (2.3,.75)--(3,.75);
  \node at (1,1.9) {横向收缩};
  \node at (1,-1) {单轴拉伸：$Y,\nu$};
 \end{scope}
\end{tikzpicture}
\figcap{剪切、均匀压缩、单轴拉伸。虚线是未变形轮廓；中图箭头都朝内，右图轴向伸长并伴随横向收缩。}
\label{ch12:fig:moduli}
\end{center}

通常实验的参考构形已经处于大气压下；在其上增加小载荷时，
可先扣除共同的大气压及原有压缩，研究增量应力与增量应变。
弹性模量的数值依材料而异，一些常见固体可达 $10^{10}\,\mathrm{Pa}$ 或更高量级。
单轴拉伸时侧面自由，因此 $\sigma_{yy}=\sigma_{zz}=0$。
记 $\varepsilon_{xx}=\varepsilon_\parallel$，
$\varepsilon_{yy}=\varepsilon_{zz}=\varepsilon_\perp$。横向本构式给出
\[
 0=\lambda_e(\varepsilon_\parallel+2\varepsilon_\perp)+2G\varepsilon_\perp,
 \qquad
 \nu=-\frac{\varepsilon_\perp}{\varepsilon_\parallel}
     =\frac{\lambda_e}{2(\lambda_e+G)}.
\]
$\nu$ 为泊松比（Poisson's ratio）。正泊松比材料受拉后横向收缩；也存在负泊松比材料。
再将 $\varepsilon_\perp=-\nu\varepsilon_\parallel$ 代回轴向应力，
\begin{equation}
 \sigma_{xx}=Y\varepsilon_\parallel,\qquad
 Y=\frac{G(3\lambda_e+2G)}{\lambda_e+G}
  =2G(1+\nu)=3K(1-2\nu).
 \label{ch12:eq:young}
\end{equation}
$Y$ 是杨氏模量（Young's modulus），$Y,G,K,\lambda_e$ 的量纲均为压强。
稳定各向同性材料要求剪切与体压缩能量均为正，即 $G>0,K>0$，从而 $-1<\nu<1/2$。
对于原长 $L$、截面积 $A$ 的均匀细杆，
$F/A=Y\Delta L/L$，故整体等效劲度系数 $k=YA/L$；
材料的刚硬程度与整个物体的劲度不是同一个物理量。

\begin{origexamplebox}{自重引起的伸长}
\label{ch12:ex:01}
长为 $L$、截面积为 $A$、密度为 $\rho$、杨氏模量为 $Y$ 的均匀细杆竖直悬挂，
求自重引起的总伸长。取小形变，$L,A,\rho$ 均指参考构形的量。
\end{origexamplebox}
\begin{xesolution}
从杆的自由下端向上量原构形坐标 $x$。图\ref{ch12:fig:selfweight}右侧取出截面以下的部分：
它长 $x$、质量 $\rho Ax$，上方截面的拉力平衡其全部重量，因此
\[
 F(x)=\rho Agx,\qquad \varepsilon(x)=\frac{F(x)}{YA}=\frac{\rho gx}{Y}.
\]
\begin{center}
\begin{tikzpicture}[>=Latex,scale=.85]
 \draw[thick] (-.9,3.5)--(.9,3.5);
 \draw[thick] (-.25,0) rectangle(.25,3.5);
 \draw[guide] (-.6,1.7)--(.7,1.7);
 \draw[<->] (-.7,0)--(-.7,1.7) node[midway,left] {$x$};
 \draw[<->] (.8,0)--(.8,3.5) node[midway,right] {$L$};
 \node[below] at (0,0) {自由端};
 \begin{scope}[shift={(4,0)}]
  \draw[thick] (-.25,0) rectangle(.25,1.7);
  \draw[vector] (0,1.7)--(0,2.7) node[above] {$F(x)$};
  \draw[vector] (0,.85)--(0,-.5) node[below] {$\rho Axg$};
  \node[right] at (.5,1) {截面以下的杆};
 \end{scope}
\end{tikzpicture}
\figcap{自重拉伸的坐标与受力。右图的向上拉力作用于截面，向下重力作用于所取部分。}
\label{ch12:fig:selfweight}
\end{center}
微段原长 $\dd x$，伸长为 $\varepsilon(x)\dd x$，所以
\[
 \boxed{\Delta L=\int_0^L\frac{\rho gx}{Y}\dd x=\frac{\rho gL^2}{2Y}}.
\]
$F$ 从下端的零增至上端的 $\rho AgL$，故平均应变恰为最大值的一半。
截面积消去，是因为增大截面同时增大承重和抗拉能力。$\rho gL/Y$ 无量纲，
$\Delta L$ 是长度；小形变还要求 $\rho gL/Y\ll1$。非均匀杆须重算每一截面以下的重量。
\end{xesolution}

\section{细杆纵波与弹性能}

设侧面自由的均匀细杆截面积为 $A$，纵向位移为 $u(x,t)$。
$x$ 标记原构形位置；图\ref{ch12:fig:rod-element}中的两端实际位置分别为
$x+u(x,t)$、$x+\dd x+u(x+\dd x,t)$，因此
\[
 \dd L-\dd x=u(x+\dd x,t)-u(x,t)=u_x\dd x,\qquad\varepsilon=u_x.
\]
用下标 $x,t$ 表示偏导数。微元右端受力向右为 $A\sigma(x+\dd x)$，
左端受力向左为 $A\sigma(x)$，其差为 $A\sigma_x\dd x$，而质量为 $\rho A\dd x$。
于是
\begin{equation}
 \rho A\dd x\,u_{tt}=A\sigma_x\dd x,\qquad
 u_{tt}=\frac{Y}{\rho}u_{xx},\qquad c_{\rm rod}=\sqrt{\frac{Y}{\rho}}.
 \label{ch12:eq:rod-wave}
\end{equation}
这里已用 $\sigma=Yu_x$。该一维杆波近似还要求波长远大于横向尺寸，
横向惯性可以忽略，侧面能够随泊松效应自由伸缩。

\begin{center}
\begin{tikzpicture}[>=Latex,scale=.88]
 \draw[thick] (-.5,0)--(8,0);
 \draw[thick] (1,-.2)--(1,.2) (4,-.2)--(4,.2);
 \node[above] at (1,.2) {$x$}; \node[above] at (4,.2) {$x+\dd x$};
 \draw[<->] (1,-.5)--(4,-.5) node[midway,below] {$\dd x$};
 \draw[thick,DeepBlue] (-.5,-2)--(8,-2);
 \draw[thick,DeepBlue] (1.5,-2.2)--(1.5,-1.8) (5,-2.2)--(5,-1.8);
 \draw[->,TextGray] (1,-1.3)--(1.5,-1.3) node[midway,above] {$u(x)$};
 \draw[->,TextGray] (4,-1.3)--(5,-1.3) node[midway,above] {$u(x+\dd x)$};
 \node[below] at (3.25,-2.2) {$\dd L=\dd x+u_x\dd x$};
 \draw[thick] (1.5,-4.3) rectangle(5,-3.6);
 \draw[vector] (1.5,-3.95)--(-.3,-3.95) node[above] {$A\sigma(x)$};
 \draw[vector] (5,-3.95)--(7.2,-3.95) node[above] {$A\sigma(x+\dd x)$};
 \node[below] at (3.25,-4.4) {$\dd m=\rho A\dd x$};
\end{tikzpicture}
\figcap{细杆微元的三层图：原位置、变形后位置、取出微元后的两端受力。位移差决定应变，应力差决定加速度。}
\label{ch12:fig:rod-element}
\end{center}

单位参考体积的动能密度为 $\rho u_t^2/2$。准静态地把应变从零增加到 $\varepsilon$，
外力每单位体积所做的功为
$\int_0^\varepsilon \sigma(\eta)\dd\eta=\int_0^\varepsilon Y\eta\dd\eta=Y\varepsilon^2/2$。
因此总能量密度为
\begin{equation}
 w=\frac12\rho u_t^2+\frac12Yu_x^2.
 \label{ch12:eq:rod-energy}
\end{equation}
$\sigma\,\dd\varepsilon$ 量纲是 $\mathrm{J/m^3}$；不能把它当作整根杆的能量。

\section{弦的横波、两个偏振与能流}

张力为 $T$、线密度为 $\mu$ 的柔弦沿 $x$ 轴伸展，两个横向位移分别为
$\xi(x,t),\eta(x,t)$，空间位置为 $(x,\xi,\eta)$。
相邻点的切向线元与单位切向量为
\[
 \dd\vect r=(1,\xi_x,\eta_x)\dd x,\qquad
 \vect t=\frac{(1,\xi_x,\eta_x)}{\sqrt{1+\xi_x^2+\eta_x^2}}
 \simeq(1,\xi_x,\eta_x).
\]
小斜率 $|\xi_x|,|\eta_x|\ll1$ 时，张力大小到一阶视为恒定。
图\ref{ch12:fig:string-element}中微元受力为右端 $T\vect t(x+\dd x)$ 与左端 $-T\vect t(x)$；
两者的横向差给出
\begin{equation}
 \mu\,\xi_{tt}=T\xi_{xx},\qquad
 \mu\,\eta_{tt}=T\eta_{xx},\qquad c_{\rm string}=\sqrt{\frac T\mu}.
 \label{ch12:eq:string-wave}
\end{equation}
这里 $T$ 是力，单位 $\mathrm N$；$\mu$ 是质量除以长度，故 $T/\mu$ 的量纲为速度平方。
两个横向自由度具有相同波速，不表示弦元沿传播方向随波前进。

\begin{center}
\begin{tikzpicture}[x={(1cm,0cm)},y={(-.5cm,-.42cm)},z={(0cm,1cm)},>=Latex]
 \draw[axis] (0,0,0)--(6.4,0,0) node[right] {$x$};
 \draw[axis] (0,0,0)--(0,2.2,0) node[below left] {$y\ (\xi)$};
 \draw[axis] (0,0,0)--(0,0,1.9) node[above] {$z\ (\eta)$};
 \draw[curve,domain=.25:5.8,samples=90,variable=\q]
  plot ({\q},{.65*sin(50*\q)},{.5*cos(50*\q)});
 \coordinate(P) at(3,{.65*sin(150)},{.5*cos(150)});
 \coordinate(Q) at(3.7,{.65*sin(185)},{.5*cos(185)});
 \fill (P) circle(1.7pt); \fill(Q) circle(1.7pt);
 \draw[vector] (P)--++(-1.1,{-.65*.872665*cos(150)*1.1},{.5*.872665*sin(150)*1.1});
 \draw[vector] (Q)--++(1.1,{.65*.872665*cos(185)*1.1},{-.5*.872665*sin(185)*1.1});
 \node[above] at (2.25,0,.35) {$-T\vect t(x)$};
 \node[below] at (5,0,-1.1) {$T\vect t(x+\dd x)$};
 \draw[guide] (P)--(3,0,0); \draw[guide] (Q)--(3.7,0,0);
 \draw[<->] (3,0,1.1)--(3.7,0,1.1) node[midway,above] {$\dd x$};
\end{tikzpicture}
\figcap{右手坐标系中的弦微元。实线由同一空间曲线生成，两个力箭头沿其端点切线、背离所取微元。}
\label{ch12:fig:string-element}
\end{center}

弦元长度
$\dd s=\sqrt{1+\xi_x^2+\eta_x^2}\dd x
\simeq[1+(\xi_x^2+\eta_x^2)/2]\dd x$。
在恒定预张力近似下，多出的势能为 $T(\dd s-\dd x)$，所以单位长度的能量为
\begin{equation}
 \mathcal E=\frac12\mu(\xi_t^2+\eta_t^2)+\frac12T(\xi_x^2+\eta_x^2).
 \label{ch12:eq:string-energy}
\end{equation}
以一个偏振 $y$ 为例，截面左侧通过张力对右侧做功的功率为
$P=-Ty_xy_t$。右行波 $y=A\cos(kx-\omega t)$ 满足
$y_x=-kA\sin(kx-\omega t)$、$y_t=\omega A\sin(kx-\omega t)$，
因此 $P=Tk\omega A^2\sin^2(kx-\omega t)>0$。周期平均得到
\begin{equation}
 \overline P=\frac12Tk\omega A^2
 =\frac12\mu c_{\rm string}\omega^2A^2,\qquad
 Z=\frac T{c_{\rm string}}=\mu c_{\rm string}=\sqrt{\mu T}.
 \label{ch12:eq:string-power}
\end{equation}
$Z$ 为弦的特征阻抗，单位为 $\mathrm{kg/s}$。线性、无色散近似下，
波速由介质与预应力决定，不由振幅和频率决定；固定频率下能流与振幅平方成正比。

\section{一维波动方程、行波与驻波}

对 $u_{tt}=c^2u_{xx}$，令特征坐标 $\zeta=x-ct,\ \chi=x+ct$。
链式法则给出
$\partial_x=\partial_\zeta+\partial_\chi$、
$\partial_t=-c\partial_\zeta+c\partial_\chi$，因此
\[
 u_{tt}-c^2u_{xx}=-4c^2u_{\zeta\chi}=0.
\]
先对 $\chi$ 积分，$u_\zeta$ 只依赖 $\zeta$；再对 $\zeta$ 积分，得到
\begin{equation}
 u(x,t)=f(x-ct)+g(x+ct).
 \label{ch12:eq:dalembert}
\end{equation}
第一项保持形状向右平移，第二项向左平移。若给定
$u(x,0)=F(x)$、$u_t(x,0)=V(x)$，则
$f+g=F$、$-cf'+cg'=V$，解这两式并积分可得
\[
 u(x,t)=\frac{F(x-ct)+F(x+ct)}2+
 \frac1{2c}\int_{x-ct}^{x+ct}V(s)\dd s.
\]
这说明求波的运动还需初始位移与初始速度，只有波动方程本身不足以唯一确定波形。

简谐右行波为 $u=A\cos(kx-\omega t+\varphi)$，
$\omega=ck,\ k=2\pi/\lambda,\ f=\omega/(2\pi),\ T_{\rm per}=1/f$。
$kx-\omega t+\varphi$ 为相位。同一相位点满足 $k\dd x-\omega\dd t=0$，
故速度为 $c=\omega/k$。图\ref{ch12:fig:travel-standing}左侧的瞬时零位移点会移动，
不能叫作驻波的“波节”。

两列等幅、同频、反向行波可组成
$u=2A\sin kx\cos\omega t$。其空间振幅为 $|2A\sin kx|$，
$x=n\pi/k$ 始终不动，是波节；两节之间振幅最大的地方是波腹。
固定端要求 $u=0$，反射位移反相；自由端要求横向力 $Tu_x=0$，反射位移同相。
两端固定、间距 $L$ 时 $\sin kL=0$，因此
\begin{equation}
 k_n=\frac{n\pi}{L},\quad \lambda_n=\frac{2L}{n},\quad
 f_n=\frac{nc}{2L},\qquad n=1,2,\ldots.
 \label{ch12:eq:fixed-modes}
\end{equation}

\begin{center}
\begin{tikzpicture}[>=Latex,x=.85cm,y=.75cm]
 \draw[axis] (-.3,0)--(6.6,0) node[right] {$x$};
 \draw[curve,domain=0:6,samples=100] plot(\x,{sin(60*\x)});
 \draw[guide,domain=0:6,samples=100] plot(\x,{sin(60*(\x-.65))});
 \draw[vector] (2.6,1.5)--(4,1.5) node[midway,above] {$c$};
 \node[above] at (1.5,1.05) {波峰};
 \node[below] at (4.5,-1.05) {波谷};
 \draw[<->] (1.5,-1.8)--(4.5,-1.8) node[midway,below] {$\lambda/2$};
 \node at (3,-2.8) {行波：实、虚线为不同时刻};
 \begin{scope}[shift={(9,0)}]
  \draw[axis] (-.3,0)--(6.6,0) node[right] {$x$};
  \draw[curve,domain=0:6,samples=100] plot(\x,{sin(60*\x)});
  \draw[guide,domain=0:6,samples=100] plot(\x,{-.5*sin(60*\x)});
  \foreach \p in {0,3,6}{\fill(\p,0)circle(2pt);}
  \node[above] at (1.5,1.05) {波腹};
  \draw[thin,TextGray] (3,-.15)--(3,-1.4);
  \node[below] at (3,-1.55) {波节};
  \node at (3,-2.8) {驻波：波节位置不随时间改变};
 \end{scope}
\end{tikzpicture}
\figcap{行波与驻波的区别。左侧虚线是平移后的波形；右侧虚线仅改变各点的振幅比例，波节保持不动。}
\label{ch12:fig:travel-standing}
\end{center}

\begin{origexamplebox}{弦的基频与张力}
\label{ch12:ex:02}
长 $L$、线密度 $\mu$ 的弦两端固定，基频为 $f_1$。
若长度、线密度不变而基频提高为 $2f_1$，张力应变为原来的多少倍？
\end{origexamplebox}
\begin{xesolution}
基频对应两端各一个波节、中间一个波腹，由式\eqref{ch12:eq:fixed-modes}，
$f_1=c/(2L)=\sqrt{T/\mu}/(2L)$。两次状态相除，
\[
 \frac{2f_1}{f_1}=\frac{\sqrt{T'/\mu}}{\sqrt{T/\mu}}
 =\sqrt{\frac{T'}T},\qquad \boxed{T'=4T}.
\]
频率加倍先要求波速加倍，而张力与波速平方成正比，所以不是两倍。
若张力变化使弦的长度、线密度明显改变，就必须同时计入这些变化。
\end{xesolution}

两个横向分量可以有不同相位。令
$\xi=A\cos(kx-\omega t)$、
$\eta=B\cos(kx-\omega t+\delta)$。
在固定材料点处，$\delta=0,\pi$ 时运动轨迹为直线，称线偏振；
$A=B,\delta=\pm\pi/2$ 时轨迹为圆，称圆偏振；一般情形为椭圆偏振。
这是介质点在垂直传播方向平面内的运动轨迹，不是波的传播轨迹。
\subsection{波动方程不等于完整的运动答案}
质点运动必须给初位置、初速度；连续介质则要对每一个位置给初位移、初速度。波动方程只规定变化规律，初始条件才决定是哪一次运动，边界条件又决定传播到端点后的反射方式。
对无界弦，右行项 $f(x-ct)$ 在 $x$ 增加 $c\Delta t$、时间增加 $\Delta t$ 后保持相同的自变量，所以波形向右；左行项 $g(x+ct)$ 同理向左。固定材料点只随时间上下振动，不能把波形的水平平移理解成弦本身向右滑动。

\begin{origexamplebox}{松手后分裂成两份的弦脉冲}\label{ch12:ex:teaching-initial-pulse}
一根足够长的均匀弦，波速为 $c$。在 $t=0$ 时拉成位移形状 $F(x)$ 并从静止释放，即 $y(x,0)=F(x)$、$y_t(x,0)=0$；$F$ 为局域的平滑函数。求后续波形，并解释为什么两列波的振幅各为原来的一半，能量却没有丢失。
\end{origexamplebox}
\begin{xesolution}
写 $y=f(x-ct)+g(x+ct)$。初位移给 $f(x)+g(x)=F(x)$；初速度给 $-cf'(x)+cg'(x)=0$。后式说明 $f-g$ 为常量。局域条件使远处位移趋零，常量可取零，故
\[
 y(x,t)=\frac12F(x-ct)+\frac12F(x+ct).
\]
每一份以速度 $c$ 向相反方向传播。代回 $t=0$，两半重叠恢复 $F$，两者的速度贡献反号，确实使初速度为零。
设弦张力为 $\mathcal T$、线密度为 $\mu$，满足 $\mathcal T=\mu c^2$。初始全是弹性势能：
\[
 E_0=\frac{\mathcal T}{2}\int_{-\infty}^{\infty}[F'(x)]^2\dd x.
\]
对右行半幅脉冲，$y_x=F'/2$、$y_t=-cF'/2$，动能与势能相等，其总能量为
\[
 E_+=\int\left(\frac{\mu c^2}{8}[F']^2+
 \frac{\mathcal T}{8}[F']^2\right)\dd x=\frac12E_0.
\]
左行脉冲同样带走 $E_0/2$。把“一半振幅对应四分之一能量”直接套在初始静止弦形上，会漏掉行波中的动能；初始状态与单向行波的动能条件不同。该解还要求端点反射尚未返回所研究区域。
\end{xesolution}

\section{张紧薄膜的二维波动（选读）}

膜的两个空间自变量为 $x,y$，横向位移只有一个分量 $z(x,y,t)$。
设参考面密度为 $\mu_s$，各向同性的单位长度张力为 $\sigma_s$。
$\sigma_s$ 的单位为 $\mathrm{N/m}$，不同于弦张力 $T$ 的 $\mathrm N$，
也不同于三维应力的 $\mathrm{N/m^2}$。这一模型取预张力恒定、小斜率，忽略弯曲刚度。

取膜上一块区域 $D$，从 $+z$ 方向看，其边界按逆时针方向绕行。
图\ref{ch12:fig:membrane}中的边界线元与向上单位法向到一阶为
\[
 \dd\vect r=(\dd x,\dd y,z_x\dd x+z_y\dd y),\qquad
 \vect n\simeq(-z_x,-z_y,1).
\]
周围膜对这块膜的力沿膜面、垂直边界，并指向区域外，所以
\[
 \dd\vect F=\sigma_s\,\dd\vect r\times\vect n
 =\sigma_s(\dd y,-\dd x,z_x\dd y-z_y\dd x)+O(|\nabla z|^2).
\]
检验方向：在区域右边界，边界线元沿 $+y$，而 $\dd\vect r\times\vect n$ 的水平分量沿 $+x$，
确实朝外，不能把叉积次序颠倒。环绕闭曲线时 $\oint\dd x=\oint\dd y=0$，
故一阶水平合力消失。竖直合力用格林公式变成
\begin{align}
 F_z&=\sigma_s\oint_{\partial D}(z_x\dd y-z_y\dd x)\\
 &=\sigma_s\iint_D(z_{xx}+z_{yy})\dd x\dd y.
 \label{ch12:eq:membrane-force}
\end{align}
这里所用公式为 $\oint(P\dd x+Q\dd y)=\iint(Q_x-P_y)\dd x\dd y$，
取 $P=-z_y,Q=z_x$。将任意区域的质量乘以加速度，与合力相等，就得到局部方程
\begin{equation}
 z_{tt}=\frac{\sigma_s}{\mu_s}(z_{xx}+z_{yy}),\qquad
 c_{\rm membrane}=\sqrt{\frac{\sigma_s}{\mu_s}}.
 \label{ch12:eq:membrane-wave}
\end{equation}

\begin{center}
\begin{tikzpicture}[x={(1cm,0cm)},y={(-.45cm,-.80cm)},z={(0cm,1cm)},>=Latex]
 \draw[axis] (0,0,0)--(5.1,0,0) node[right] {$x$};
 \draw[axis] (0,0,0)--(0,4.3,0) node[below left] {$y$};
 \draw[axis] (0,0,0)--(0,0,2.4) node[above] {$z$};
 \foreach \q in {0,1,2,3}{
  \draw[TextGray!45,domain=0:4,samples=45,variable=\s]
    plot({\s},{\q},{.12*\s+.07*\q+.045*\s*\q});
 }
 \foreach \q in {0,1,2,3,4}{
  \draw[TextGray!45,domain=0:3,samples=40,variable=\s]
    plot({\q},{\s},{.12*\q+.07*\s+.045*\q*\s});
 }
 \draw[DeepBlue,thick] (1,1,.235)--(3,1,.565)--(3,2,.77)--(1,2,.35)--cycle;
 \draw[vector] (3,1.25,.61625)--(3,1.8,.729);
 \coordinate(P) at(3,1.5,.6675);
 \draw[vector] (P)--++({1.25*(1+.205*.205)},{-1.25*.1875*.205},{1.25*.1875});
 \draw[vector,TextGray] (P)--++(-.1875,-.205,1);
 \node[right=4pt] at (4.4,1.5,.93) {$\dd\vect F$};
 \node[above=4pt] at (2.8125,1.295,1.75) {$\vect n$};
 \draw[thin,TextGray] (3.04,1.78,.733)--(3.9,3.2,.50);
 \node[below right=3pt] at (3.9,3.2,.50) {$\dd\vect r$};
 \draw[thin,TextGray] (1.5,1.55,.393125)--(.1,2.8,.45);
 \node[left=4pt] at (.1,2.8,.45) {$D$};
\end{tikzpicture}
\figcap{右手坐标系中的膜面及有向边界。曲面网格、边界、切线、法向和外向膜力均由同一曲面 $z=0.12x+0.07y+0.045xy$ 生成。}
\label{ch12:fig:membrane}
\end{center}

面积元由两个切向量的叉积求得：
$\dd S=\sqrt{1+z_x^2+z_y^2}\dd x\dd y$。
在小斜率近似下，面积增加为
$\dd S-\dd x\dd y\simeq(z_x^2+z_y^2)\dd x\dd y/2$。
恒定单位长度张力所对应的势能增加等于 $\sigma_s$ 乘以面积增加，
故单位参考面积的能量为
\begin{equation}
 \mathcal E_s=\frac12\mu_s z_t^2+\frac12\sigma_s(z_x^2+z_y^2).
 \label{ch12:eq:membrane-energy}
\end{equation}
这里常数平面面积的能量可以选为零，并不意味着可以省略形变面积的变化。

无限均匀膜上的平面简谐波可写为
$z=\operatorname{Re}\{A e^{\ii(\omega t-k_xx-k_yy)}\}$，
代入波动方程给出
$\omega^2=c_{\rm membrane}^2(k_x^2+k_y^2)$。
频率决定 $|\vect k|$，$k_x:k_y$ 及其符号决定传播方向，
复数 $A$ 的模和幅角分别决定振幅与初相位；这些是不同的自由度。

矩形膜的边长为 $a,b$，四边固定时试取
$z=X(x)Y(y)\cos(\omega t+\phi)$。边界要求 $X(0)=X(a)=Y(0)=Y(b)=0$，
因此 $X\propto\sin(m\pi x/a)$、$Y\propto\sin(n\pi y/b)$，代回方程得到
\begin{equation}
 z_{mn}\propto\sin\frac{m\pi x}{a}\sin\frac{n\pi y}{b}\cos(\omega_{mn}t+\phi),
 \quad
 \omega_{mn}=\pi\sqrt{\frac{\sigma_s}{\mu_s}}
 \sqrt{\frac{m^2}{a^2}+\frac{n^2}{b^2}}.
 \label{ch12:eq:membrane-modes}
\end{equation}
$m,n$ 均为正整数。不同的空间方向各带一个模态指标，若 $a=b$，
交换 $m,n$ 不改变频率，可能出现简并。圆膜需用极坐标和贝塞尔函数
（Bessel functions）；允许频率由形状与边界共同决定，而不是只由材料决定。

\section{三维弹性体与纵横体波（选读）}

\subsection{从应力散度到运动方程}

对任意物质体积 $D$，外向面力合力的第 $i$ 分量为
\[
 F_i=\oint_{\partial D}\sigma_{ij}n_j\dd S
 =\iiint_D\frac{\partial\sigma_{ij}}{\partial x_j}\dd V.
\]
第二步对每个分量使用高斯定理。若体积力密度为 $f_i$，单位 $\mathrm{N/m^3}$，
牛顿第二定律给出
\begin{equation}
 \rho u_{i,tt}=\partial_j\sigma_{ij}+f_i.
 \label{ch12:eq:stress-div}
\end{equation}
小形变时，在这一阶精度内可将参考坐标导数与空间坐标导数等同。
均匀各向同性介质中，代入式\eqref{ch12:eq:lame}，逐项求导：
\[
 \partial_j\sigma_{ij}
 =\lambda_e\partial_i(\partial_k u_k)
   +G\partial_j(\partial_i u_j+\partial_j u_i)
 =(\lambda_e+G)\partial_i(\nabla\cdot\vect u)+G\nabla^2u_i.
\]
无体积力时得到纳维—柯西方程（Navier--Cauchy equation）
\begin{equation}
 \rho\vect u_{tt}=(\lambda_e+G)\nabla(\nabla\cdot\vect u)+G\nabla^2\vect u
 =(\lambda_e+2G)\nabla(\nabla\cdot\vect u)-G\nabla\times(\nabla\times\vect u).
 \label{ch12:eq:navier}
\end{equation}
末式用到 $\nabla^2\vect u=\nabla(\nabla\cdot\vect u)-\nabla\times(\nabla\times\vect u)$。
为看清耦合，第一分量完全展开为
\[
 \rho u_{x,tt}=(\lambda_e+2G)u_{x,xx}
 +G(u_{x,yy}+u_{x,zz})
 +(\lambda_e+G)(u_{y,xy}+u_{z,xz});
\]
另外两式由 $x,y,z$ 循环替换得到。横向二阶导数不能从展开式中漏去。

\subsection{三维能量密度}

从零应变沿直线路径增加到 $\boldsymbol\varepsilon$，应力随路径参数也线性增加，
所以单位体积弹性能为
\begin{equation}
 w_{\rm el}=\frac12\sum_{i,j}\sigma_{ij}\varepsilon_{ij}
 =\frac{\lambda_e}{2}(\operatorname{tr}\boldsymbol\varepsilon)^2
   +G\sum_{i,j}\varepsilon_{ij}^2.
 \label{ch12:eq:elastic-energy}
\end{equation}
双重求和中的 $\varepsilon_{xy}^2$ 与 $\varepsilon_{yx}^2$ 都要计入；
纯剪切时便得到 $w_{\rm el}=G\gamma_{xy}^2/2$，与 $\tau=G\gamma_{xy}$ 的做功一致。
加上动能密度 $\rho|\vect u_t|^2/2$ 就得到总机械能密度。
把应变分成体积部分和无迹部分，还可写成
\[
 w_{\rm el}=\frac K2(\operatorname{tr}\boldsymbol\varepsilon)^2
 +G\left|\boldsymbol\varepsilon-\frac13
 (\operatorname{tr}\boldsymbol\varepsilon)\mathsf I\right|^2,
\]
清楚显示 $K>0,G>0$ 是稳定性条件。

\subsection{平面波的三种偏振与不同波速}

取平面波 $\vect u=\vect A e^{\ii(\vect k\cdot\vect r-\omega t)}$ 的实部。
代入式\eqref{ch12:eq:navier}得到
\[
 \rho\omega^2\vect A=(\lambda_e+G)\vect k(\vect k\cdot\vect A)+Gk^2\vect A.
\]
将振幅唯一分成
$\vect A_P=(\vect A\cdot\hat{\vect k})\hat{\vect k}$ 与
$\vect A_S=\vect A-\vect A_P$。
前者与 $\vect k$ 平行，满足 $\vect k\times\vect A_P=0$；
后者与 $\vect k$ 垂直，满足 $\vect k\cdot\vect A_S=0$。
分别投影原式：
\begin{equation}
 c_P^2=\frac{\omega_P^2}{k^2}=\frac{\lambda_e+2G}{\rho}
      =\frac{K+4G/3}{\rho},\qquad
 c_S^2=\frac{\omega_S^2}{k^2}=\frac G\rho.
 \label{ch12:eq:body-speeds}
\end{equation}
一个传播方向对应一种纵偏振、两种互相独立的横偏振，如图\ref{ch12:fig:polarization}。
对一般波场，在适当边界条件下可作无旋与无散的赫尔姆霍兹分解；
对运动方程分别取散度与旋度，也得到
$(\nabla\cdot\vect u)_{tt}=c_P^2\nabla^2(\nabla\cdot\vect u)$、
$(\nabla\times\vect u)_{tt}=c_S^2\nabla^2(\nabla\times\vect u)$。
因此上述结论不只是某一个特殊方向的平面波巧合。

\begin{center}
\begin{tikzpicture}[>=Latex]
 \foreach \q/\name in {0/纵偏振,4.5/横偏振一,9/横偏振二}{
 \begin{scope}[shift={(\q,0)}]
  \draw[axis] (0,0)--(2.7,0) node[right] {$\vect k$};
  \node at (1.2,-1) {\name};
 \end{scope}}
 \draw[<->,DeepBlue,very thick] (.6,.5)--(1.9,.5) node[midway,above] {$\vect A_P$};
 \draw[<->,DeepBlue,very thick] (5.7,-.6)--(5.7,.9) node[above] {$\vect A_{S1}$};
 \draw[DeepBlue,thick] (10.2,.4) circle(.16);
 \fill[DeepBlue] (10.2,.4) circle(1.6pt);
 \node[above] at (10.2,.75) {$\vect A_{S2}$：出入纸面};
\end{tikzpicture}
\figcap{沿同一方向传播的三种独立体波偏振。箭头与点圆符号表示材料点振动方向，不是另外的传播方向。}
\label{ch12:fig:polarization}
\end{center}

细杆纵波的 $\sqrt{Y/\rho}$ 不能替换为体纵波的 $c_P$：
前者侧面可以横向收缩，后者平面纵波伴随的横向约束使有效纵向模量为 $\lambda_e+2G$。
稳定各向同性材料满足 $c_P>c_S$。地震中先到达的压缩体波称P波（primary wave），
较慢的剪切体波称S波（secondary wave）。液态外核不能传递剪切体波；
这处边界称古登堡不连续面（Gutenberg discontinuity），核幔边界深约 $2900\,\mathrm{km}$，波的折射、反射和模式转换为认识地球内部提供依据。
S波在边界可转换为能穿过液体的压缩波，不能将“液体不传S波”说成“所有S波能量均原样反射”。

\section{流体中的声波与绝热声速}

静态流体不能维持剪切应力，在无黏理想流体模型中 $G=0$，只剩压缩纵波。
设平衡密度为 $\rho_0$，小纵位移为 $u(x,t)$。质量守恒给出
$\delta\rho=-\rho_0 u_x$：伸长区域密度减小。
小压强变化为
$\delta p=(\partial p/\partial\rho)_{\rm process}\delta\rho=-K u_x$。
两端压强差形成力 $-\partial_x\delta p$，所以
\[
 \rho_0u_{tt}=-\partial_x\delta p=K u_{xx},\qquad
 c^2=\frac K{\rho_0}=\left(\frac{\partial p}{\partial\rho}\right)_{\rm process}.
\]
这里“过程”必须说明，否则 $K$ 不唯一。普通空气中声波的压缩、稀疏较快，
在一周期内、一个波长尺度上来不及通过热传导抹平温度差；忽略热耗散时取局部可逆绝热过程。
它不是说气体温度不变，而是每一物质微元与邻域的净换热可忽略。

理想气体满足 $pV^\gamma=\text{常量}$，其中 $\gamma=C_p/C_V$。
对固定质量微元求微分：
$\dd p/p=-\gamma\,\dd V/V$，于是
\begin{equation}
 K_s=-V\left(\frac{\partial p}{\partial V}\right)_s=\gamma p,\qquad
 c=\sqrt{\frac{\gamma p}{\rho_0}}
  =\sqrt{\frac{\gamma RT_0}{M_{\rm mol}}}.
 \label{ch12:eq:sound-speed}
\end{equation}
最后一步用 $pV=nRT_0$ 与 $\rho_0=nM_{\rm mol}/V$，
$M_{\rm mol}$ 为摩尔质量而不是单个分子质量。若使用单分子质量 $m_0$，
则应写 $c=\sqrt{\gamma k_{\rm B}T_0/m_0}$。

牛顿（Newton）采用等温压缩关系，会得到 $K_T=p$；
拉普拉斯（Laplace）的绝热修正使预测声速乘以 $\sqrt\gamma$。
例如 $T_0=293\,\mathrm K$、$M_{\rm mol}=0.029\,\mathrm{kg/mol}$、$\gamma=1.4$ 时，
等温模型给出约 $290\,\mathrm{m/s}$，绝热模型给出约 $343\,\mathrm{m/s}$。
这是在同一温度、同一气体参数下的比较，不能拿不同条件的历史近似数值混作精确测量。

